\chapter{Le dépôt, cahier de laboratoire}

Un texte mathématique change sans cesse : une hypothèse ajoutée, une notation corrigée, une
démonstration réécrite. Trois mois plus tard on veut savoir quand l'hypothèse est apparue,
et pourquoi. Un dossier de fichiers numérotés \fichier{article-v3-final-bis.tex} ne le dit
pas. Un gestionnaire de versions le dit, à condition qu'on lui ait donné de quoi répondre.

\section{Le dépôt et ses commits}

L'outil est \texttt{git}. Un dépôt est un dossier dont \texttt{git} garde l'histoire : chaque
état enregistré, appelé \emph{commit}, porte la liste des fichiers modifiés, la différence
ligne à ligne, une date, un auteur et un message. On peut revenir à n'importe quel état,
comparer deux états, et chercher dans l'histoire.

\begin{terminal}
# créer un dépôt dans le dossier courant
git init
# inscrire une modification
git add article.tex
git commit -m "Lemme 3 : la finitude est nécessaire (cf. 3.4)"
# l'histoire, un commit par ligne
git log --oneline
\end{terminal}

Le dépôt local se double d'une copie sur un service d'hébergement (GitHub, GitLab, ou un
serveur de l'université). La copie distante sert de sauvegarde, de point de rencontre avec
les collaborateurs et, on le verra, de déclencheur pour les vérifications automatiques.

\section{Un commit par idée}

\begin{regle}
Un commit contient une seule idée, et son message dit pourquoi.
\end{regle}

« Corrige le lemme 3 » dit quoi. « Lemme 3 : l'hypothèse de finitude est nécessaire,
contre-exemple en 3.4 » dit pourquoi, et c'est cela qu'on cherchera. Le message est écrit
pour un lecteur qui n'a pas le contexte : soi-même dans six mois, ou une collaboratrice.

Avec une idée par commit, on défait une idée sans toucher aux autres. Si la correction du
lemme 3 et le renommage d'une notation sont dans le même commit, on ne peut annuler l'un
sans l'autre.

\begin{exemple}
Dans le projet qui sert d'exemple, un commit typique dit en substance : « Problème 14 : la règle
des signes le long de chaque droite suffit, avec un triangle coupé en deux quand aucune
quatrième droite n'évite ses sommets ; preuve en Lean, vérification sur 37\,213 cas ». Une phrase
suffit à savoir ce qui a changé et où regarder.
\end{exemple}

\section{Retrouver quand et pourquoi}

Deux commandes répondent à presque toutes les questions d'histoire.

\begin{terminal}
# les commits qui ont ajouté ou retiré ce mot
git log -S "finitude" -- article.tex
# pour chaque ligne, le dernier commit qui l'a touchée
git blame article.tex
\end{terminal}

La première retrouve le moment où une hypothèse est apparue. La seconde montre, ligne à
ligne, de quel commit vient le texte actuel ; le message de ce commit dit pourquoi.

\section{Figer une version}

Quand un article part chez un éditeur, on veut pouvoir retrouver exactement le texte
envoyé, même après des mois de corrections. Une \emph{étiquette} (tag) nomme un commit.

\begin{terminal}
# nommer l'état actuel
git tag article-soumis-2026-10
git push origin article-soumis-2026-10
# ce qui a changé depuis l'envoi
git diff article-soumis-2026-10 -- article.tex
\end{terminal}

La dernière commande donne la liste des changements à joindre à une version révisée.

\section{Ne jamais réécrire l'histoire publiée}

\begin{warningblock}
Une fois un commit envoyé sur la copie distante, on ne le modifie plus. Les commandes qui
réécrivent l'histoire (\cmd{git push --force}, \cmd{git rebase} sur une branche partagée)
détruisent le travail des autres et la trace qu'on voulait garder.
\end{warningblock}

Une erreur se corrige par un nouveau commit qui la répare et le dit. L'erreur reste visible
dans l'histoire, à côté de sa correction. Tant mieux : on sait qu'elle a existé, et l'on
peut chercher ce qu'elle a entraîné entre-temps.

\section{Branches, ou pas}

Les équipes de développement travaillent sur des \emph{branches} qu'elles fusionnent
ensuite. Pour un projet mené seul, une seule branche principale suffit souvent. Les
branches deviennent utiles quand on contribue au dépôt d'un autre (on y reviendra avec
mathlib au chapitre~\ref{ch:lean}) ou quand deux chantiers longs doivent avancer sans se
gêner.
